%% ===================================================================================
%% Capítulo 5 -- Projeto da Arquitetura e da Solução
%% FONTE DA VERDADE da arquitetura, do contrato de API, do roteamento e da estratégia de
%% cache. Os diagramas citados estão em doc/diagrams/src/ (ver doc/diagrams/README.md).
%% Dono: Arquiteto/Redator. Mudanças propostas pelos desenvolvedores entram pelo mural.
%% Convenção: "normativo" = contrato a ser cumprido; "esboço" = ilustração não vinculante.
%% ===================================================================================
\chapter{Projeto da Arquitetura e da Solução}\label{cap:arquitetura}

Este capítulo define \emph{como} o Vitrine atende aos requisitos do Capítulo~\ref{cap:requisitos}. Suas seções são normativas para a implementação, exceto os trechos explicitamente chamados de \emph{esboço}. Quando a implementação precisar divergir, o desenvolvedor propõe a mudança no mural e o Arquiteto atualiza este capítulo antes da integração.

\section{Visão Arquitetural e Princípios}\label{sec:arq-visao}

A arquitetura é um \textbf{monólito modular}: um único repositório, um único processo Node.js e um único arquivo de banco de dados. Os principais direcionadores são a portabilidade e a simplicidade de implantação (\req{RNF01}), o desempenho de leitura (\req{RNF02}), a segurança (\req{RNF03}) e a manutenibilidade por domínios (\req{RNF08}). Esses direcionadores se traduzem nos princípios a seguir.

\begin{enumerate}[label=\textbf{PA\arabic*.}, leftmargin=1.6cm]
  \item \textbf{Monólito modular.} Um processo, um banco, um repositório; módulos separados por domínio (Dados, Core, UI), com dependências unidirecionais (Figura~\ref{fig:arq-camadas}).
  \item \textbf{Pré-renderizar por padrão.} Páginas públicas são estáticas e regeneradas de forma incremental; o dinamismo fica restrito ao painel administrativo e à API.
  \item \textbf{Acesso a dados centralizado.} Somente \arq{src/lib/db.js} instancia o cliente do Prisma; rotas e componentes não escrevem consultas.
  \item \textbf{Segurança em profundidade.} Verificação otimista no \en{Proxy}, verificação de sessão em cada Route Handler e Server Component protegido, e validação de entrada no servidor.
  \item \textbf{Configuração por ambiente.} Segredos e parâmetros de implantação em variáveis de ambiente, nunca no repositório \cite{twelvefactor}.
  \item \textbf{Conteúdo é dado, não código.} O corpo é Markdown; HTML bruto é proibido (RN11).
  \item \textbf{Personalização por dados.} Identidade e cores vêm do banco e chegam ao CSS por variáveis, sem novo \en{build}.
  \item \textbf{Simplicidade operacional.} Sem fila, cache externo ou serviço adicional obrigatório; a implantação é um contêiner Docker com um volume para o banco (\adr{ADR-011}).
\end{enumerate}

\section{Decisões Arquiteturais}\label{sec:arq-decisoes}

As decisões estruturantes estão registradas como ADR no Apêndice~\ref{ap:adr}. A Tabela~\ref{tab:resumo-adr} as resume. As ADR-011 a ADR-014 foram acrescentadas em 30/09/2026, após o merge M-001 e a atualização P-001, as ADR-015 a ADR-017 em 01/10/2026, quando o Engenheiro Chefe encerrou as decisões que restavam, as ADR-018 e ADR-019 em 01/10/2026, após o merge M-004, e a ADR-020 em 01/10/2026, após o merge M-005; os identificadores D-xx das decisões iniciais continuam citados no texto e levam à ADR que as registra.

\begingroup\small\renewcommand{\arraystretch}{1.25}
\begin{longtable}{@{}L{2.0cm}L{5.6cm}L{6.4cm}@{}}
\caption{Resumo das decisões arquiteturais aceitas}\label{tab:resumo-adr}\\
\toprule
\textbf{ADR} & \textbf{Decisão} & \textbf{Motivo principal}\\
\midrule
\endfirsthead
\toprule
\textbf{ADR} & \textbf{Decisão} & \textbf{Motivo principal}\\
\midrule
\endhead
\bottomrule
\endfoot
\adr{ADR-001} & Monólito modular com Next.js (App Router) & Simplicidade operacional; uma unidade de implantação\\
\adr{ADR-002} & SQLite como banco embutido & Portabilidade; arquivo único; custo zero\\
\adr{ADR-003} & Prisma ORM 7 como camada de acesso & Modelo declarativo, migrações e cliente tipado\\
\adr{ADR-004} & better-auth no lugar do Auth.js previsto na IC & E-mail e senha com sessão em banco e adaptador Prisma\\
\adr{ADR-005} & ISR com revalidação sob demanda & Desempenho estático com atualização imediata\\
\adr{ADR-006} & Markdown com remark/rehype e KaTeX & Escrita científica nativa, sem HTML bruto\\
\adr{ADR-007} & Tematização por variáveis CSS e Tailwind 4 & Personalização sem \en{build}\\
\adr{ADR-008} & JavaScript com JSDoc, sem TypeScript & Herdado do projeto base; menor fricção\\
\adr{ADR-009} & Documentation-Driven Development com \LaTeX{} e diagramas como código & Coerência entre agentes\\
\adr{ADR-010} & Fluxo Git por domínios com mural na \texttt{main} & Coordenação assíncrona rastreável\\
\adr{ADR-011} & Implantação por imagem Docker em VPS tradicional, com volume para o SQLite & Persistência do arquivo de banco e ambiente reprodutível (resolve D-04)\\
\adr{ADR-012} & Estado de publicação \texttt{DRAFT}/\texttt{PUBLISHED}, padrão \texttt{PUBLISHED} & Rascunhos sem histórico de versões (resolve D-15)\\
\adr{ADR-013} & Seis campos opcionais de redes sociais em \texttt{Configuration} & Rodapé configurável sem tabela adicional (resolve D-19)\\
\adr{ADR-014} & Adaptador \texttt{better-sqlite3} e conformidade com o Prisma 7 & O Prisma 7 exige \en{driver adapter} para o SQLite (resolve D-09)\\
\adr{ADR-015} & Nome Vitrine e licença MIT & Identidade e licença permissiva do código aberto (resolve D-01 e D-02)\\
\adr{ADR-016} & Documento de Arquitetura e Requisitos de Software, sem formato acadêmico & Documentação a serviço do produto (resolve D-18)\\
\adr{ADR-017} & Decisões técnicas com valor padrão adotado & Encerra as questões técnicas iniciais, refináveis por RFC (D-03, D-05 a D-08, D-10 a D-14, D-16, D-17)\\
\adr{ADR-018} & Administrador por \texttt{npm run admin:create} e contrato da API (envelope \texttt{data}, erros de entrada e de origem) & Aprova as nove propostas do Core (RFC P-002; fecha o D-03)\\
\adr{ADR-019} & Manter o Prisma 7 apesar de quatro alertas do \texttt{npm audit} & A correção automática rebaixaria o Prisma para a versão 6\\
\adr{ADR-020} & Texto sobre a cor primária escolhido por \texttt{readableTextOn} e indicador de foco neutro & O branco fixo reprova o nível AA com a paleta padrão, e o foco não pode depender de uma cor configurável\\
\end{longtable}
\fonte{Elaborado pelos autores.}
\endgroup

\section{Visões da Arquitetura}\label{sec:arq-visoes}

\subsection{Contexto}

O sistema atende dois atores humanos --- Leitor e Administrador --- e depende apenas de uma VPS tradicional com Docker (\decisao{D-04}). O Implantador obtém o código de um repositório Git, constrói a imagem Docker e a executa na VPS (Figura~\ref{fig:arq-contexto}).

\vitrinefigura[tipo=diagrama, arquivo=arq-contexto, agente=Arquiteto, largura=0.9,
  origem={doc/diagrams/src/arq-contexto.puml},
  descricao={Diagrama de contexto (C4 nível 1): Leitor, Administrador, Implantador, o sistema Vitrine, a VPS tradicional com Docker e o repositório Git.}]
  {Diagrama de contexto do Vitrine}{fig:arq-contexto}

\subsection{Contêineres}

Na terminologia do modelo C4, há um único \emph{contêiner} de aplicação --- o servidor Next.js, que na implantação é empacotado em uma imagem Docker (\adr{ADR-011}) --- mais o navegador do usuário e o arquivo SQLite. Dentro do servidor convivem as páginas públicas (Server Components com ISR), o painel administrativo, os Route Handlers, o better-auth, o Proxy, a camada de dados e o cache de ISR (Figura~\ref{fig:arq-conteineres}).

\vitrinefigura[tipo=diagrama, arquivo=arq-conteineres, agente=Arquiteto, largura=0.95,
  origem={doc/diagrams/src/arq-conteineres.puml},
  descricao={Diagrama de contêineres (C4 nível 2): navegador, servidor Node.js em contêiner Docker com páginas públicas, painel, Route Handlers, better-auth, Proxy, camada de dados (Prisma e adaptador better-sqlite3) e cache ISR, e o banco SQLite em volume.}]
  {Diagrama de contêineres do Vitrine}{fig:arq-conteineres}

\subsection{Componentes e Regras de Dependência}

A Figura~\ref{fig:arq-componentes} mostra os componentes por diretório e o agente responsável por cada um, e a Figura~\ref{fig:arq-camadas} resume as camadas. As regras de dependência são normativas:

\begin{itemize}
  \item componentes em \arq{src/components} \textbf{não} importam \arq{src/lib/db.js} nem o cliente gerado do Prisma (\arq{@/generated/prisma}); recebem dados por \en{props} dos Server Components ou os buscam nas rotas \arq{/api};
  \item rotas em \arq{src/app/api} \textbf{não} escrevem consultas Prisma; chamam funções de \arq{db.js};
  \item apenas \arq{db.js} instancia o cliente do Prisma na aplicação, com o adaptador \arq{better-sqlite3} (uma instância por processo, \adr{ADR-014}); a única exceção é o \en{script} \arq{prisma/seed.js}, executado pela CLI fora do Next.js;
  \item toda rota de escrita valida a sessão e a entrada antes de chamar a camada de dados.
\end{itemize}

\vitrinefigura[tipo=diagrama, arquivo=arq-componentes, agente=Arquiteto, largura=0.95,
  origem={doc/diagrams/src/arq-componentes.puml},
  descricao={Diagrama de componentes por diretório (src/app, src/components, src/lib, prisma), coloridos por agente responsável, com as regras de dependência.}]
  {Componentes do monólito e regras de dependência}{fig:arq-componentes}

\vitrinefigura[tipo=diagrama, arquivo=arq-camadas, agente=Arquiteto, largura=0.85,
  origem={doc/diagrams/src/arq-camadas.puml},
  descricao={Camadas Apresentação, Aplicação, Dados e Infraestrutura, com o sentido permitido das dependências.}]
  {Camadas e sentido das dependências}{fig:arq-camadas}

\section{Estrutura do Repositório e Convenções}\label{sec:arq-estrutura}

O Código~\ref{lst:arvore-atual} mostra a estrutura verificada no repositório na data desta redação, após os merges M-001 a M-005: o \en{scaffold} do Next.js, o esquema do Prisma 7 com a migração inicial e o \en{seed}, as funções de dados, a autenticação, as rotas de API de conteúdo e de configuração, o comando do administrador, o tema e os primeiros componentes da UI (\arq{src/components/admin/} e \arq{src/components/ui/}), os testes e a governança. O Código~\ref{lst:arvore-alvo} mostra a estrutura-alvo, que os agentes constroem por módulo, já com as pastas criadas pelos agentes fora do desenho inicial: \arq{scripts/}, \arq{tests/} (\arq{setup/}, \arq{dados/}, \arq{core/} e \arq{ui/}), \arq{src/app/api/_lib/} e \arq{src/app/admin/_lib/}. Os diretórios \arq{_lib} começam por sublinhado e, no App Router do Next.js, são \emph{pastas privadas}: não geram rota e guardam o código compartilhado pelos Route Handlers e pelas páginas da mesma área. O alias \arq{@/} aponta para \arq{src/} (definido em \arq{jsconfig.json}).

\begin{lstlisting}[caption={Estrutura atual do repositório (verificada)}, label={lst:arvore-atual}]
.
|-- CLAUDE.md  ARQUITETO.md  DESENVOLVEDORES.md  MURAL_DE_AGENTES.md  LICENSE
|-- doc/                      # documentação (ref/, tex/, diagrams/)
|-- prisma/                   # schema.prisma, migrations/ (init) e seed.js (Dados)
|-- prisma.config.ts          # configuração da CLI (datasource.url vem de DATABASE_URL)
|-- scripts/create-admin.js   # npm run admin:create (Core)
|-- tests/                    # setup/ (comum), dados/, core/ e ui/ (Vitest)
|-- vitest.config.mjs         # configuração do Vitest (alias @, banco temporário)
|-- .env.example              # nomes das variáveis de ambiente, sem valores
|-- public/                   # ativos estáticos do scaffold
|-- src/proxy.js              # verificação otimista de sessão em /admin (Core)
|-- src/app/                  # layout.js e page.js (scaffold); globals.css (tema, UI)
|   |-- login/                # page.js (cartão de login; usa o LoginForm)
|   |-- admin/                # layout.js, page.js e _lib/ (session, actions)
|   `-- api/                  # auth/[...all], contents/, configuration/ e _lib/
|-- src/components/           # admin/ (painel, LoginForm) e ui/ (Button, icons)
|-- src/lib/                  # db.js, content-tree.js (Dados); auth.js, auth-client.js,
|                             #   validators.js, revalidate.js (Core); theme.js (UI)
|-- src/generated/prisma/     # cliente gerado; ignorado pelo Git (ADR-014)
|-- eslint.config.mjs  jsconfig.json  next.config.mjs  postcss.config.mjs
|-- package.json  package-lock.json  skills-lock.json
`-- .agents/ .claude/ .windsurf/   # skills do Prisma para os agentes
\end{lstlisting}

\begin{lstlisting}[caption={Estrutura-alvo do código (normativa)}, label={lst:arvore-alvo}]
src/
|-- proxy.js                      # verificação otimista de sessão (Core)
|-- app/
|   |-- layout.js                 # RootLayout: fontes, lang="pt-BR", tema (Core+UI)
|   |-- globals.css               # Tailwind 4 + tokens de tema (UI)
|   |-- page.js                   # página inicial
|   |-- [...slug]/page.js         # leitura de documento/coleção (ISR)
|   |-- not-found.js
|   |-- login/page.js             # cartão de login; usa o LoginForm (UI)
|   |-- admin/                    # painel (dinâmico): layout, page, conteudo/, configuracoes/
|   |   `-- _lib/                 # pasta privada (Core)
|   |       |-- session.js        # requireAdminSession
|   |       `-- actions.js        # Server Function signOut ("use server")
|   `-- api/
|       |-- _lib/                 # http.js (envelope, erros, sessão, origem) e contents.js
|       |-- auth/[...all]/route.js
|       |-- contents/route.js
|       |-- contents/[id]/route.js
|       |-- contents/reorder/route.js
|       `-- configuration/route.js
|-- components/                   # layout/, markdown/, editor/, admin/, ui/ (UI)
`-- lib/
    |-- db.js                     # único ponto de acesso ao Prisma (Dados)
    |-- content-tree.js           # montagem da árvore e resolução de caminhos (Dados)
    |-- auth.js  auth-client.js   # better-auth servidor e cliente (Core)
    |-- validators.js             # esquemas de entrada (Core)
    |-- revalidate.js             # regras de revalidação (Core)
    `-- theme.js                  # variáveis CSS e contraste (UI)
src/generated/prisma/             # cliente gerado por "prisma generate" (não versionado)
prisma/                           # schema.prisma, migrations/, seed.js (Dados)
prisma.config.ts                  # CLI do Prisma: schema, migrations, seed, DATABASE_URL
scripts/create-admin.js           # npm run admin:create: administrador inicial (Core)
tests/                            # setup/ (comum), dados/, core/ e ui/ (Vitest)
vitest.config.mjs                 # configuração dos testes
Dockerfile  docker-compose.yml  .dockerignore   # implantação (M13)
.env.example                      # nomes das variáveis de ambiente, sem valores
\end{lstlisting}

Convenções: identificadores em inglês; comentários e textos de interface em português; módulos ES (\texttt{import}/\texttt{export}); funções públicas de \arq{src/lib} documentadas com JSDoc (\adr{ADR-008}); o arquivo \arq{.env.example} lista os nomes das variáveis de ambiente, sem valores (o \arq{.gitignore} ignora \arq{.env*} e traz a exceção \arq{!.env.example}). Código compartilhado por Route Handlers ou páginas de uma mesma área fica em uma pasta privada \arq{_lib} dessa área, e não em \arq{src/lib}, que é reservado aos módulos de uso geral. Scripts de linha de comando ficam em \arq{scripts/} e são executados com o \texttt{tsx}, pois os módulos de \arq{src/} usam o alias \arq{@/} e o cliente gerado do Prisma é TypeScript. O cliente do Prisma é gerado em \arq{src/generated/prisma}, diretório ignorado pelo Git, e importado de \arq{@/generated/prisma/client} somente por \arq{db.js}; por isso \texttt{prisma generate} deve rodar depois de \texttt{npm install} e antes de \texttt{next build} (\adr{ADR-014}).

\section{Modelo de Dados}\label{sec:arq-dados}

\subsection{Modelo de Domínio}

O domínio possui seis entidades (Figura~\ref{fig:classes-dominio}): quatro de autenticação, gerenciadas pelo better-auth (\texttt{User}, \texttt{Session}, \texttt{Account}, \texttt{Verification}), e duas de negócio do Vitrine (\texttt{Content} e \texttt{Configuration}), além do enumerado \texttt{ContentStatus}. \texttt{Content} possui um auto-relacionamento que forma a árvore de documentos e coleções, um estado de publicação (\texttt{DRAFT} ou \texttt{PUBLISHED}, padrão \texttt{PUBLISHED}; \adr{ADR-012}) e se relaciona com o autor (\texttt{User}). \texttt{Configuration} é um registro único e guarda também os seis links sociais do rodapé, todos opcionais (\adr{ADR-013}).

\vitrinefigura[tipo=diagrama, arquivo=classes-dominio, agente=Arquiteto, largura=0.95,
  origem={doc/diagrams/src/classes-dominio.puml},
  descricao={Diagrama de classes do domínio espelhando prisma/schema.prisma: User, Session, Account, Verification, Content (auto-relacionamento pai e filhos, com status), Configuration (com os seis links sociais), o enumerado ContentStatus e o valor lógico ContentType.}]
  {Modelo de domínio do Vitrine}{fig:classes-dominio}

\subsection{Modelo Físico}

O modelo físico (Figura~\ref{fig:er-vitrine}) é implementado por \arq{prisma/schema.prisma}. O dicionário de dados completo está no Apêndice~\ref{ap:dicionario} e o esquema no Apêndice~\ref{ap:schema}. As convenções adotadas são: identificadores \texttt{cuid}; nomes de tabela pela diretiva \texttt{@@map} (\texttt{user}, \texttt{session}, \texttt{account}, \texttt{verification}, \texttt{contents}, \arq{configurations}); nomes de coluna iguais aos campos do modelo (\emph{camelCase}); carimbos \texttt{createdAt} e \texttt{updatedAt}; e exclusão em cascata nas relações, exceto o autor do conteúdo (\texttt{Content.author}), que usa \texttt{Restrict}: excluir um usuário que tenha conteúdo é bloqueado (RNF10). O enumerado \texttt{ContentStatus} é gravado como texto; o SQLite não impõe o conjunto de valores, de modo que a verificação cabe ao Prisma e à aplicação. Os seis links sociais são colunas de texto opcionais da tabela \arq{configurations}: \arq{discordUrl}, \arq{githubUrl}, \arq{youtubeUrl}, \arq{instagramUrl}, \arq{linkedinUrl} e \arq{twitchUrl}.

\vitrinefigura[tipo=diagrama, arquivo=er-vitrine, agente=Arquiteto, largura=0.95,
  origem={doc/diagrams/src/er-vitrine.dbml},
  descricao={Diagrama entidade-relacionamento gerado no dbdiagram.io a partir de er-vitrine.dbml: tabelas user, session, account, verification, contents (auto-relacionamento e coluna status do enumerado ContentStatus) e configurations (com as seis colunas de links sociais).}]
  {Diagrama entidade-relacionamento do Vitrine}{fig:er-vitrine}

A Figura~\ref{fig:ic-er} reproduz o modelo do relatório de IC, com três tabelas. A evolução consiste em substituir a tabela genérica \texttt{users} pelas tabelas do better-auth e em explicitar as relações de cascata.

\vitrinefigura[tipo=ic, arquivo=ic-fig2-er, largura=0.35,
  fonte={Relatório de IC \cite{reis2026ic}, Figura 2.}]
  {Diagrama entidade-relacionamento original, do relatório de IC}{fig:ic-er}

\subsection{Hierarquia de Conteúdo}\label{sec:arq-hierarquia}

A árvore de conteúdos é representada pelo campo \texttt{parentId} (nulo para nós raiz). A URL pública é a cadeia de \en{slugs} da raiz ao nó (RN04). Para a navegação, a árvore é carregada por \textbf{uma única consulta}, sem o corpo dos documentos, e montada em memória em tempo linear; a resolução de um caminho percorre os níveis comparando \en{slugs}. O Código~\ref{lst:algoritmo-arvore} é um \emph{esboço} do algoritmo a ser implementado em \arq{src/lib/content-tree.js}. O parâmetro \texttt{onlyPublished} implementa a RN15: um nó \texttt{DRAFT} é removido junto com toda a sua subárvore, de modo que um descendente publicado de um ancestral em rascunho também não aparece. A área pública usa a árvore podada; o painel usa a árvore completa.

\begin{lstlisting}[language=VJS, caption={Esboço: montagem da árvore e resolução de caminho}, label={lst:algoritmo-arvore}]
// rows: [{ id, title, slug, type, status, order, parentId, createdAt }] (sem o corpo)
export function buildTree(rows, { onlyPublished = false } = {}) {
  const byId = new Map(rows.map((r) => [r.id, { ...r, children: [] }]));
  const roots = [];
  for (const node of byId.values()) {
    const parent = node.parentId ? byId.get(node.parentId) : null;
    (parent ? parent.children : roots).push(node);
  }
  // RN15: nó não publicado oculta toda a sua subárvore
  const prune = (list) =>
    list
      .filter((n) => !onlyPublished || n.status === "PUBLISHED")
      .map((n) => ({ ...n, children: prune(n.children) }));
  const sortLevel = (list) => {
    list.sort((a, b) => a.order - b.order || a.createdAt - b.createdAt); // RN05
    list.forEach((n) => sortLevel(n.children));
  };
  const visible = prune(roots);
  sortLevel(visible);
  return visible;
}

export function resolvePath(roots, slugs) {
  let level = roots;
  let node = null;
  for (const slug of slugs) {
    node = level.find((n) => n.slug === slug) ?? null;
    if (!node) return null; // 404 (RF14, RN15)
    level = node.children;
  }
  return node;
}
\end{lstlisting}

A alternativa de gravar o caminho materializado em cada nó foi descartada nesta versão, pois exigiria atualizar todos os descendentes a cada renomeação ou movimentação; fica registrada como evolução possível, para o caso de o volume de conteúdo o justificar (\decisao{D-06}, \adr{ADR-017}).

\subsection{Observações sobre o Esquema Atual}\label{sec:arq-observacoes-schema}

A revisão do esquema (\arq{prisma/schema.prisma}) contra esta especificação resultou nas observações da Tabela~\ref{tab:obs-schema}, que foram entradas para o agente Dados. O merge M-001, aprovado em 30/09/2026, integrou a conformidade com o Prisma 7 (\adr{ADR-014}), o enumerado \texttt{ContentStatus} com a coluna \texttt{status} (\adr{ADR-012}) e os seis links sociais de \texttt{Configuration} (\adr{ADR-013}). O merge M-002, aprovado em 01/10/2026, resolveu as observações 1, 3, 4 e 5 na migração inicial e manteve a 2 como validação na aplicação; o merge M-003 aplicou a observação 8 na camada de dados (filtro \texttt{PUBLISHED} nas leituras públicas, Seção~\ref{sec:arq-visibilidade}).

\begin{longtable}{@{}L{0.8cm}L{5.6cm}L{2.2cm}L{4.8cm}@{}}
\caption{Observações sobre o esquema atual do Prisma}\label{tab:obs-schema}\\
\toprule
\textbf{\#} & \textbf{Observação} & \textbf{Regra ou RNF} & \textbf{Recomendação}\\
\midrule
\endfirsthead
\toprule
\textbf{\#} & \textbf{Observação} & \textbf{Regra ou RNF} & \textbf{Recomendação}\\
\midrule
\endhead
\bottomrule
\endfoot
1 & \texttt{slug} sem restrição de unicidade & RN03 & \emph{Resolvido (M-002)}: índice único (\texttt{parentId}, \texttt{slug}) mais índice parcial para a raiz, em SQL (Seção~\ref{sec:arq-indice-parcial}; D-05)\\[1pt]
2 & \texttt{type} é texto livre, ao passo que \texttt{status} já é um enumerado & RN02 & \emph{Mantido como texto}: DOCUMENT e COLLECTION são validados na aplicação (Zod, Core), pois o SQLite não impõe enumerados e um \texttt{enum} do Prisma só duplicaria a regra\\[1pt]
3 & Sem índice em \texttt{parentId} & RNF02 & \emph{Resolvido (M-002)}: índice (\texttt{parentId}, \texttt{order}) para a consulta da árvore\\[1pt]
4 & \arq{Content.author} com \texttt{onDelete: Cascade}: excluir o usuário apaga todo o seu conteúdo & RNF10 & \emph{Resolvido (M-002)}: \texttt{onDelete: Restrict}; excluir um usuário que tenha conteúdo é bloqueado pelo banco\\[1pt]
5 & \texttt{Configuration} sem garantia de registro único & RN08 & \emph{Resolvido (M-002)}: identificador fixo \texttt{default} (padrão do esquema) e acesso por \emph{upsert}\\[1pt]
6 & Datas e booleanos sem valor padrão em \arq{User.emailVerified} & --- & Manter: o valor é fornecido pelo better-auth\\[1pt]
7 & Tabela \texttt{verification} sem relação com \texttt{user} & --- & Manter: desenho do better-auth\\[1pt]
8 & \texttt{status} tem padrão \texttt{PUBLISHED}: um conteúdo recém-criado torna-se público de imediato (decisão D-15) & RN15, RN16, risco R11 & Manter o padrão; a API deve aceitar a criação já como \texttt{DRAFT}, a interface deve exibir o estado com destaque e a camada de dados deve filtrar por \texttt{PUBLISHED} nas leituras públicas\\[1pt]
\end{longtable}
\fonte{Elaborado pelos autores.}

\subsection{Índice Parcial da Raiz}\label{sec:arq-indice-parcial}

A regra RN03 exige \en{slug} único entre irmãos. O índice único composto (\texttt{parentId}, \texttt{slug}) resolve a regra para todos os nós com pai, mas não para a raiz: no SQLite, \texttt{NULL} é distinto de \texttt{NULL}, de modo que dois conteúdos de raiz (\texttt{parentId} nulo) com o mesmo \en{slug} passariam pelo índice composto. O agente Dados resolveu a lacuna com um \textbf{índice único parcial} escrito à mão no SQL da migração inicial (Código~\ref{lst:indice-parcial}), que cobre apenas as linhas da raiz. O Prisma não modela índices parciais, mas também não os trata como divergência (\en{drift}): \texttt{migrate diff} contra as migrações e contra o banco não acusa diferenças. O resultado é a garantia de unicidade no próprio banco, sem depender só da validação na aplicação, que permanece como primeira barreira (\decisao{D-05}, \adr{ADR-017}).

\begin{lstlisting}[caption={Índice parcial da raiz (trecho da migração inicial)}, label={lst:indice-parcial}]
CREATE UNIQUE INDEX "contents_parentId_slug_key"
  ON "contents"("parentId", "slug");

-- manual: cobre somente a raiz
CREATE UNIQUE INDEX "contents_root_slug_key"
  ON "contents"("slug") WHERE "parentId" IS NULL;
\end{lstlisting}

Como a migração é SQL versionado, qualquer nova migração gerada pelo Prisma deve preservar esse índice; o agente Dados deve conferir isso ao alterar o esquema. Em banco descartável, o agente Dados verificou que um \en{slug} repetido na raiz e entre irmãos é recusado e que o mesmo \en{slug} em pais diferentes é aceito; os testes de integração correspondentes (casos CT08 e CT28) ainda devem ser automatizados.

\subsection{Migrações e Dados Iniciais}

As alterações de esquema são versionadas por migrações do Prisma (\arq{prisma/migrations}), geradas em desenvolvimento e aplicadas em produção por um comando não interativo (Apêndice~\ref{ap:implantacao}). O \en{seed} insere a configuração padrão (RF23) e cria, sob controle, o primeiro administrador (RF22). A primeira migração (\arq{20261001034824_init}), o \en{seed} idempotente da configuração padrão (\arq{prisma/seed.js}, executado por \texttt{tsx} porque o cliente gerado é TypeScript) e a instância do cliente em \arq{src/lib/db.js} foram integrados pelo merge M-002; as funções de acesso a dados foram integradas pelo merge M-003; o administrador inicial não passa pelo \en{seed}: é criado pelo comando \texttt{npm run admin:create} do Core, integrado no merge M-004 (\decisao{D-03}, \adr{ADR-018}). Como \arq{prisma} e \arq{dotenv} estão em \texttt{dependencies}, a imagem de produção aplica as migrações com \texttt{prisma migrate deploy}. No Prisma 7, o comando \texttt{migrate dev} não gera o cliente nem executa o \en{seed} automaticamente, e o \en{seed} é declarado em \arq{prisma.config.ts} (chave \texttt{migrations.seed}), não mais no \arq{package.json}; por isso \texttt{prisma generate} e \texttt{prisma db seed} são passos explícitos. Um cuidado adicional: um caminho relativo em \arq{DATABASE_URL} pode ser resolvido de forma diferente pela CLI e pelo adaptador em tempo de execução, criando silenciosamente um segundo banco; o agente Dados verificou que, com \texttt{file:./dev.db} e execução a partir da raiz, a CLI e o adaptador abrem o mesmo arquivo; o adaptador resolve o caminho a partir do diretório de trabalho do processo, e em produção o caminho deve ser absoluto (por exemplo, \texttt{file:/data/vitrine.db}). As instruções de CLI da versão 7 estão nas \en{skills} do Prisma instaladas em \arq{.claude/skills/}.

\section{Autenticação e Autorização}\label{sec:arq-auth}

\subsection{Configuração}

O better-auth 1.7.6 é configurado em \arq{src/lib/auth.js} (merge M-004) com o adaptador do Prisma para SQLite, autenticação por e-mail e senha, sessão persistida no banco, cadastro público desabilitado e o \en{plugin} de integração de \en{cookies} para o Next.js. As rotas são expostas por um único Route Handler em \arq{src/app/api/auth/[...all]/route.js}. A sessão no servidor é obtida por \texttt{getSession(headers)}, usada pelos Route Handlers e pelas páginas do painel. O cliente React (\arq{src/lib/auth-client.js}) usa a mesma origem da página e dispensa variáveis \texttt{NEXT\_PUBLIC\_*}. Variáveis de ambiente: \arq{DATABASE_URL}, \arq{BETTER_AUTH_SECRET} e \arq{BETTER_AUTH_URL} (nomes confirmados na versão instalada, que as lê sozinha) e, lidas somente pelo comando do administrador inicial, \arq{ADMIN_EMAIL}, \arq{ADMIN_PASSWORD} e \arq{ADMIN_NAME}. A Tabela~\ref{tab:auth-config} resume a configuração efetiva, aprovada na \adr{ADR-018}.

\begingroup\small
\begin{longtable}{@{}L{3.0cm}L{9.6cm}L{1.9cm}@{}}
\caption{Configuração efetiva do better-auth}\label{tab:auth-config}\\
\toprule
\textbf{Aspecto} & \textbf{Configuração} & \textbf{Origem}\\
\midrule
\endfirsthead
\toprule
\textbf{Aspecto} & \textbf{Configuração} & \textbf{Origem}\\
\midrule
\endhead
\bottomrule
\endfoot
Cadastro & Desabilitado (\texttt{disableSignUp}); \texttt{POST /api/auth/sign-up/email} responde 400 & RF22\\
\addlinespace
Senha & De 8 a 128 caracteres; o \en{hash} é do próprio better-auth, o mesmo no login e no comando do administrador & RNF03\\
\addlinespace
Sessão & Persistida no banco; expira em 7 dias e, em uso, é renovada no máximo uma vez por dia & RN12\\
\addlinespace
\en{Cookie} & \arq{better-auth.session_token}, \texttt{HttpOnly} e \texttt{SameSite=Lax}; com \arq{BETTER_AUTH_URL} em \texttt{https}, ganha \texttt{Secure} e o prefixo \texttt{\_\_Secure-} & RNF03\\
\addlinespace
Limitação de tentativas & Ligada em todos os ambientes, com contadores em memória por endereço IP (cabeçalho \texttt{x-forwarded-for}) e por rota: 100 requisições por 60 s no geral e \textbf{5 por 60 s em \texttt{/sign-in/email}} (resposta 429) & RNF03\\
\addlinespace
Telemetria & Desligada & RNF15\\
\addlinespace
Segredo & \arq{BETTER_AUTH_SECRET} obrigatório em produção; sem ele, as rotas de autenticação e de API respondem 500 (falha fechada: nenhuma rota protegida é servida) e o \en{build} registra o erro do segredo padrão sem falhar & RNF03\\
\end{longtable}
\fonte{Elaborado pelos autores, a partir de \arq{src/lib/auth.js} e da verificação do Arquiteto no merge M-004.}
\endgroup

\begin{lstlisting}[language=VJS, caption={Configuração efetiva do servidor de autenticação (resumo de src/lib/auth.js)}, label={lst:auth-esboco}]
export const auth = betterAuth({
  appName: "Vitrine",
  database: prismaAdapter(prisma, { provider: "sqlite" }),
  emailAndPassword: {
    enabled: true,
    disableSignUp: true,            // RF22
    minPasswordLength: 8,
    maxPasswordLength: 128,
  },
  session: { expiresIn: 7 * 86400, updateAge: 86400 },   // RN12
  rateLimit: {
    enabled: true, window: 60, max: 100,
    customRules: { "/sign-in/email": { window: 60, max: 5 } },
  },
  telemetry: { enabled: false },    // RNF15
  plugins: [nextCookies()],         // deve ser o último
});
// src/app/api/auth/[...all]/route.js
export const { GET, POST } = toNextJsHandler(auth);
\end{lstlisting}

A limitação de tentativas identifica o cliente pelo cabeçalho \texttt{x-forwarded-for}, que o better-auth 1.7.6 só aceita quando contém um \emph{único} endereço (um valor com vários endereços pode ter sido forjado pelo cliente). Quando o endereço não pode ser determinado, todos os clientes passam a compartilhar um único contador por rota --- por exemplo, apenas cinco tentativas de login por minuto no total ---, o que mantém a proteção, mas deixa um atacante bloquear momentaneamente o login do administrador. Por isso o \en{proxy} reverso deve \emph{definir} esse cabeçalho com o endereço real do cliente, substituindo o valor recebido, e não acrescentar-lhe um segundo endereço (Apêndice~\ref{ap:implantacao}).

\subsection{Fluxo e Proteção em Camadas}

O fluxo de login e de acesso à área administrativa está na Figura~\ref{fig:seq-autenticacao}. A proteção é feita em camadas, porque o Proxy é, por documentação do Next.js, uma verificação otimista e não uma solução completa de autorização \cite{nextjsdocs}:

\begin{table}[htb]
\caption{Camadas de proteção das rotas administrativas}\label{tab:camadas-protecao}
\small
\begin{tabularx}{\linewidth}{@{}L{3.3cm}L{5.2cm}Y@{}}
\toprule
\textbf{Camada} & \textbf{Mecanismo} & \textbf{Finalidade}\\
\midrule
1. Proxy & \arq{src/proxy.js} (casador \arq{/admin/:path*}) lê o \en{cookie} de sessão, sem consultar o banco, e redireciona para \arq{/login?redirectTo=...} se ausente & Melhor experiência: evita renderizar a página para quem não tem sessão\\
\addlinespace
2. Layout/página do painel & Server Component valida a sessão no servidor e redireciona se inválida; a página repete a validação, pois layouts não são renderizados de novo a cada navegação & Barreira real para páginas\\
\addlinespace
3. Route Handlers & Cada handler valida a sessão no banco antes de qualquer ação (RF21) e, nas mutações, antes disso, a origem da requisição; só então lê o corpo & Barreira real para a API\\
\addlinespace
4. Validação de entrada & Esquemas de entrada rejeitam dados inválidos antes da camada de dados & Integridade e segurança\\
\bottomrule
\end{tabularx}
\fonte{Elaborado pelos autores.}
\end{table}

\textbf{Interface do login e do encerramento da sessão.} O formulário de login é o componente \texttt{LoginForm} (\arq{src/components/admin/LoginForm.js}, Componente Cliente), que chama \texttt{authClient.signIn.email}: em qualquer erro de credencial mostra uma única mensagem genérica (UC08, fluxo A1); no status \texttt{429} avisa que as tentativas foram excedidas; sem resposta do servidor, informa a falha de conexão; e, em caso de sucesso, faz a navegação completa para o destino que o servidor já validou (\texttt{safeRedirectPath}). O encerramento da sessão (RF20) \emph{não} passa pela rota \arq{/api/auth/sign-out}: o botão ``Sair'' do cabeçalho do painel envia um formulário à \en{Server Function} \texttt{signOut}, definida em \arq{src/app/admin/_lib/actions.js} (diretiva \texttt{"use server"}), que chama \texttt{auth.api.signOut} com os cabeçalhos da requisição --- apaga a sessão do banco e remove o \en{cookie} --- e redireciona para \arq{/login}. O layout do painel, que é um arquivo de rota do Core, entrega a função ao cabeçalho por \en{prop}; assim, os componentes da UI não importam o módulo de autenticação do servidor (regras de dependência, Seção~\ref{sec:arq-visoes}). Toda função exportada de um arquivo \texttt{"use server"} é um ponto de entrada HTTP, alcançável por uma requisição direta; por isso esse arquivo só pode conter funções que validem no servidor tudo o que exigem (a \texttt{signOut} não recebe argumentos).

\vitrinefigura[tipo=diagrama, arquivo=seq-autenticacao, agente=Arquiteto, largura=0.95,
  origem={doc/diagrams/src/seq-autenticacao.puml},
  descricao={Sequência de autenticação: acesso a /admin sem cookie, redirecionamento a /login, POST de credenciais ao better-auth, criação de sessão, Set-Cookie e acesso autorizado.}]
  {Sequência de autenticação do administrador}{fig:seq-autenticacao}

\subsection{Política de Cadastro e Papéis}

Na versão inicial há um único perfil, o administrador. O primeiro usuário é criado pelo comando \texttt{npm run admin:create} (\decisao{D-03}, \adr{ADR-018}), executado por quem instala o sistema: ele lê \arq{ADMIN_EMAIL}, \arq{ADMIN_PASSWORD} e, opcionalmente, \arq{ADMIN_NAME}, valida o formato e chama \texttt{ensureInitialAdmin} (\arq{src/lib/auth.js}), que usa o adaptador interno do better-auth --- com o mesmo \en{hash} de senha do login --- porque a rota de cadastro está desabilitada. O comando é \textbf{idempotente}: com o mesmo e-mail, nada muda (nem a senha) e a saída é 0; recusa um segundo usuário com outro e-mail, pois há um único administrador (\decisao{D-13}; saída 1); completa a credencial de um usuário que ficou sem ela por uma execução interrompida; e nunca imprime a senha. Não há página aberta de primeiro acesso, para evitar o sequestro da instalação. As credenciais não passam pelo \arq{prisma/seed.js}, que cria apenas a configuração padrão (RF23); o comando roda com o \texttt{tsx}, como o \en{seed}. Na imagem Docker: \texttt{docker compose run -{}-rm vitrine npm run admin:create} (Apêndice~\ref{ap:implantacao}). O \en{plugin} \texttt{admin} do better-auth, presente no pacote instalado, é o caminho previsto para papéis e permissões futuros (\decisao{D-13}).

\section{Renderização, Cache e ISR}\label{sec:arq-isr}

\subsection{Estratégia por Rota}

A Tabela~\ref{tab:rotas-render} define o modo de renderização de cada conjunto de rotas. O projeto usa a ISR \emph{sem} Cache Components (o recurso \arq{cacheComponents} não está habilitado em \arq{next.config.mjs}), conforme o guia correspondente da documentação do Next.js 16 \cite{nextjsdocs}.

\begin{table}[htb]
\caption{Modo de renderização por conjunto de rotas}\label{tab:rotas-render}
\small
\begin{tabularx}{\linewidth}{@{}L{3.3cm}L{3.0cm}Y@{}}
\toprule
\textbf{Rota} & \textbf{Modo} & \textbf{Regras}\\
\midrule
\arq{/} e \arq{/[...slug]} & Estático com ISR & \arq{generateStaticParams} lista os caminhos \emph{publicados} (RN15); \texttt{revalidate} alto (por exemplo, 3600 s) como rede de segurança; novos caminhos são gerados sob demanda; revalidação sob demanda pelas mutações\\
\addlinespace
\arq{/login} & Estático ou dinâmico & Formulário; sem dados sensíveis\\
\addlinespace
\arq{/admin/**} & Dinâmico & Lê a sessão por \en{cookies}; nunca em cache\\
\addlinespace
\arq{/api/**} & Dinâmico & Route Handlers; respostas sem cache\\
\bottomrule
\end{tabularx}
\fonte{Elaborado pelos autores.}
\end{table}

\subsection{Regras de Revalidação}

A Tabela~\ref{tab:revalidacao} implementa as regras RN10 e RN16. A invalidação é feita em Route Handlers com \arq{revalidatePath}, função que só pode ser chamada no servidor, não em Componentes Cliente nem no Proxy. Em Route Handlers, a função marca o caminho para revalidação, que ocorre na próxima visita \cite{nextjsdocs}. A escolha por \emph{layout} raiz quando a estrutura muda é necessária porque a \en{sidebar} e a barra superior, presentes em todas as páginas, dependem da árvore inteira.

\begin{table}[htb]
\caption{Eventos de mutação e a revalidação correspondente}\label{tab:revalidacao}
\small
\begin{tabularx}{\linewidth}{@{}L{5.2cm}Y@{}}
\toprule
\textbf{Evento} & \textbf{Ação de revalidação}\\
\midrule
Criar ou excluir conteúdo; alterar título, \en{slug}, pai, ordem ou tipo; mover/reordenar &\texttt{revalidatePath("/", "layout")}: invalida todas as páginas, pois a árvore de navegação mudou\\
\addlinespace
Alterar o estado de publicação (\texttt{DRAFT} $\leftrightarrow$ \texttt{PUBLISHED}) de qualquer nó & \texttt{revalidatePath("/", "layout")}: a visibilidade do nó e de seus descendentes altera a árvore de navegação e o conjunto de páginas (RN16)\\
\addlinespace
Alterar somente o corpo de um documento (estado inalterado) & \texttt{revalidatePath(caminho do documento)}\\
\addlinespace
Alterar a configuração (nome, logotipo, cores, links sociais) & \texttt{revalidatePath("/", "layout")}: o rodapé aparece em todas as páginas\\
\bottomrule
\end{tabularx}
\fonte{Elaborado pelos autores, com base em RN10 e RN16.}
\end{table}

\subsection{Visibilidade por Estado de Publicação}\label{sec:arq-visibilidade}

A regra RN15 --- um conteúdo só é público se ele e todos os seus ancestrais estiverem \texttt{PUBLISHED} --- é aplicada em um único lugar, a camada de dados, e nunca nos componentes. Os pontos de aplicação são: (i) a árvore usada pela \en{sidebar}, pela barra superior e por \arq{generateStaticParams}, podada conforme o Código~\ref{lst:algoritmo-arvore}; (ii) a resolução do caminho nos Server Components, que responde 404 (\arq{notFound()}) para um nó ausente da árvore podada, de modo que um rascunho não é acessível nem por URL direta; e (iii) o mapa do site e os metadados (RF25). As rotas \arq{/api/contents} e o painel recebem a árvore completa, com o campo \texttt{status} de cada nó (RF26). Como o estado entra na regeneração do \en{layout} raiz (Tabela~\ref{tab:revalidacao}), a mudança de estado reflete-se em todas as páginas na visita seguinte. Como o padrão do estado é \texttt{PUBLISHED} (\decisao{D-15}), a proteção contra a publicação acidental depende de a interface exibir o estado de forma explícita e de a API aceitar a criação já como rascunho (risco R11).

\vitrinefigura[tipo=diagrama, arquivo=seq-leitura-isr, agente=Arquiteto, largura=0.95,
  origem={doc/diagrams/src/seq-leitura-isr.puml},
  descricao={Sequência de leitura de uma página com os três caminhos do cache: HIT, STALE e MISS, incluindo a regeneração em segundo plano e o cabeçalho x-nextjs-cache.}]
  {Sequência de leitura de uma página com ISR}{fig:seq-leitura-isr}

\vitrinefigura[tipo=diagrama, arquivo=seq-publicacao-conteudo, agente=Arquiteto, largura=0.95,
  origem={doc/diagrams/src/seq-publicacao-conteudo.puml},
  descricao={Sequência de edição e publicação: editor, PATCH na API, validação de sessão e de entrada, atualização no banco e revalidação do caminho ou do layout.}]
  {Sequência de edição e publicação de um documento}{fig:seq-publicacao-conteudo}

\subsection{Limitações e Observabilidade}

A documentação do Next.js 16 registra ressalvas que condicionam a implantação: a ISR exige o \en{runtime} Node.js e não funciona em exportação estática; com mais de uma instância, o cache padrão em sistema de arquivos é por instância e a invalidação só atinge a instância que recebeu a chamada; e o Proxy não é executado em requisições de ISR sob demanda \cite{nextjsdocs}. Por isso, a implantação usa \textbf{uma única instância}, isto é, um único contêiner (o que também convém ao SQLite). Em testes, o cabeçalho \arq{x-nextjs-cache} (\texttt{HIT}, \texttt{STALE}, \texttt{MISS}, \texttt{REVALIDATED}) permite verificar o comportamento do cache (Capítulo~\ref{cap:testes}).

\section{Pipeline de Conteúdo Markdown}\label{sec:arq-markdown}

A renderização de um documento usa o componente \arq{MarkdownRenderer} (\arq{src/components/markdown}), compartilhado pela página pública e pela pré-visualização do editor, para garantir que o autor veja o que o leitor verá. A cadeia de processamento é:

\begin{enumerate}
  \item \arq{react-markdown} converte o Markdown em elementos React, sem HTML bruto (RN11);
  \item \arq{remark-math} reconhece \texttt{\$...\$} e \texttt{\$\$...\$\$} (RF09);
  \item \arq{remark-gfm} acrescenta tabelas, tachado, listas de tarefas e links automáticos (RF11; pacote ainda não instalado, \decisao{D-08});
  \item \arq{rehype-katex} renderiza as fórmulas; a folha de estilos do KaTeX é importada uma vez no \en{layout};
  \item um \en{plugin} de realce de sintaxe colore os blocos de código (RF10; biblioteca a definir, \decisao{D-07}).
\end{enumerate}

Requisitos de segurança e de comportamento: manter a transformação padrão de URLs do \arq{react-markdown}, que restringe protocolos; não adicionar \arq{rehype-raw} sem sanitização; tratar expressões \LaTeX{} inválidas de forma tolerante (mensagem visível, sem derrubar a página); e mapear links relativos internos para as URLs por \en{slug}.

\section{Tematização Dinâmica}\label{sec:arq-tema}

A personalização visual (RF04, RF18, RF19) baseia-se em variáveis CSS alimentadas pelo banco. O mecanismo, ilustrado na Figura~\ref{fig:seq-personalizacao-tema}, tem três partes:

\begin{enumerate}
  \item \textbf{Tokens.} Em \arq{globals.css}, o bloco \texttt{@theme inline} do Tailwind mapeia os tokens de cor (por exemplo, \arq{--color-primary}) para variáveis (por exemplo, \arq{var(--vitrine-primary)}); os utilitários compilados referenciam a variável, não o valor;
  \item \textbf{Injeção.} O \en{layout} raiz, um Server Component, lê a configuração e define as variáveis no atributo \texttt{style} do elemento \texttt{<html>} --- as três cores e a cor do texto sobre a primária, \arq{--vitrine-on-primary}, calculada por \texttt{readableTextOn} (Seção~\ref{sec:arq-contraste}) ---, além do nome do site, do logotipo e do \en{favicon}; os valores são validados como \texttt{\#RRGGBB} antes de persistir, o que impede a injeção de CSS;
  \item \textbf{Invalidação.} Ao salvar a configuração, o Route Handler invalida o \en{layout} raiz; na próxima requisição, as páginas são regeneradas com as novas variáveis, sem \en{build}.
\end{enumerate}

\begin{lstlisting}[language=VCSS, caption={Tokens de tema em globals.css (resumo do arquivo integrado no merge M-005)}, label={lst:tema-css}]
@import "tailwindcss";

:root {
  /* Paleta configurável (contrato do M11): <html style> sobrescreve o padrão */
  --vitrine-primary: #3b82f6;
  --vitrine-secondary: #1e293b;
  --vitrine-accent: #f59e0b;
  /* Texto sobre a primária = readableTextOn(primária); definida no layout raiz */
  --vitrine-on-primary: #171717;
  --vitrine-on-secondary: #ffffff;

  /* Neutros fixos (D-17: sem modo escuro automático) e foco independente */
  --vitrine-background: #ffffff;  --vitrine-foreground: #171717;
  --vitrine-surface: #f5f5f5;     --vitrine-muted: #525252;
  --vitrine-border: #d4d4d4;      --vitrine-border-strong: #737373;
  --vitrine-danger: #b91c1c;
  --vitrine-focus: #171717;       --vitrine-focus-halo: #ffffff;
}

@theme inline {
  --color-primary: var(--vitrine-primary);
  --color-secondary: var(--vitrine-secondary);
  --color-accent: var(--vitrine-accent);
  --color-on-primary: var(--vitrine-on-primary);
  --color-on-secondary: var(--vitrine-on-secondary);
  /* neutros no mesmo padrão: --color-background, -foreground, -surface... */
}

/* D-17: dark: só vale sob a classe .dark, nunca pela preferência do sistema */
@custom-variant dark (&:where(.dark, .dark *));
\end{lstlisting}

O mapeamento semântico é: \textbf{primária} para ações e ênfase (botão principal e item ativo da navegação, com o fundo \texttt{bg-primary} e o texto \texttt{text-on-primary}); \textbf{secundária} para a estrutura (fundo da barra superior e cabeçalhos, com texto branco, \texttt{text-on-secondary}) e para o \en{hover} dos botões; e \textbf{destaque} para marcadores e realce não textual (marcador do item ativo e \en{hover} da navegação). O \textbf{indicador de foco} não usa a cor de destaque: é neutro e global --- contorno de 2~px em \texttt{\#171717} com halo branco (\arq{--vitrine-focus} e \arq{--vitrine-focus-halo}) ---, de modo que fica visível sobre qualquer fundo, inclusive sobre uma paleta que o administrador troque (RNF05; \adr{ADR-020}). Os demais tons da interface --- fundo, texto, superfícies, bordas e mensagens de erro --- são neutros fixos, não configuráveis. Nesta versão não há modo escuro automático: vale somente a paleta escolhida pelo administrador (\decisao{D-17}, \adr{ADR-017}), e a variante \texttt{dark:} do Tailwind só vale sob a classe \texttt{.dark}, nunca pela preferência do sistema; por isso o fundo da página é fixo e branco, o texto sobre a cor secundária é branco, o texto sobre a cor primária é escolhido por \texttt{readableTextOn} (Seção~\ref{sec:arq-contraste}) e só as três cores do tema são verificadas. Os componentes usam apenas os utilitários gerados pelos tokens (por exemplo, \texttt{bg-primary}, \texttt{text-on-primary}, \texttt{bg-surface} e \texttt{text-muted}), nunca as cores fixas do Tailwind (RF19).

\subsection{Avisos de Contraste}\label{sec:arq-contraste}

Ao salvar a configuração, o sistema calcula a razão de contraste da WCAG 2.1 (luminância relativa do sRGB, razão $(L_1 + 0{,}05)/(L_2 + 0{,}05)$) entre cada cor do tema e o branco, conforme a Tabela~\ref{tab:pares-contraste}, e emite um aviso, sem bloquear o salvamento, quando a razão é inferior ao mínimo do par (RN09, \req{RNF05}). A razão é simétrica, e o par é sempre contra o branco: o par ``primária e branco'' cobre a cor primária usada como texto sobre o fundo da página (por exemplo, em links). O texto \emph{sobre} um fundo da cor primária, ao contrário, não é branco fixo: é calculado, como descrito adiante.

\begin{table}[htb]
\caption{Pares de contraste verificados na configuração}\label{tab:pares-contraste}
\small
\begin{tabularx}{\linewidth}{@{}L{2.8cm}L{3.6cm}L{1.7cm}Y@{}}
\toprule
\textbf{Campo avisado} & \textbf{Par (cor e branco)} & \textbf{Mínimo} & \textbf{Uso que a razão protege}\\
\midrule
\texttt{primaryColor} & primária $\times$ \texttt{\#FFFFFF} & 4,5:1 & Texto na cor primária sobre a página (por exemplo, links); o texto sobre um fundo da cor primária é tratado por \texttt{readableTextOn}, e não por este par\\
\addlinespace
\texttt{secondaryColor} & secundária $\times$ \texttt{\#FFFFFF} & 4,5:1 & Texto branco sobre a barra superior e os cabeçalhos\\
\addlinespace
\texttt{accentColor} & destaque $\times$ \texttt{\#FFFFFF} & 3:1 & Marcadores e realce não textual (item ativo e \en{hover} da navegação): contraste de componentes de interface (WCAG 1.4.11); o foco não usa esta cor\\
\bottomrule
\end{tabularx}
\fonte{Elaborado pelos autores, a partir da resposta ao pedido P-002 (item 5) e da revisão do merge M-005.}
\end{table}

As funções de cálculo pertencem à UI e estão em \arq{src/lib/theme.js} (merge M-005), um módulo puro --- não importa o Next.js nem o Prisma e não usa APIs do navegador ---, que o servidor e o navegador compartilham:

\begin{itemize}
  \item \texttt{contrastRatio(a, b)} devolve a razão da WCAG 2.1 entre duas cores \texttt{\#RRGGBB}, sem arredondar e independentemente da ordem dos argumentos;
  \item \texttt{contrastWarnings(\{ primaryColor, secondaryColor, accentColor \})} devolve a lista de avisos \texttt{\{ code, field, ratio, minimum, message \}}, com o código \texttt{LOW\_CONTRAST}, vazia se nenhuma cor falha; a mensagem vem em português, por exemplo ``A cor primária tem contraste de 3,67:1 com o branco, abaixo do mínimo de 4,5:1 (WCAG 2.1 AA).'';
  \item \texttt{readableTextOn(cor)} devolve \texttt{\#FFFFFF} ou o texto escuro padrão, o de maior razão de contraste com a cor dada (no empate, o branco).
\end{itemize}

A comparação com o mínimo usa a razão \emph{sem arredondar}; o campo \texttt{ratio} e a mensagem trazem a razão \emph{truncada} em duas casas, para que 4,499 não apareça como 4,50 e pareça cumprir o mínimo (por isso a primária padrão aparece como 3,67, e não 3,68). O formulário de tema (M11) usa \texttt{contrastWarnings} na pré-visualização local, e o Route Handler de \texttt{PUT /api/configuration} (Core) a chama, depois da escrita, com as cores já gravadas --- inclusive quando o corpo da requisição não traz cores, pois o \texttt{PUT} é parcial --- e devolve a lista no campo \texttt{warnings} da resposta, ao lado de \texttt{data} (Seção~\ref{sec:arq-api}). O campo está sempre presente na resposta do \texttt{PUT} (lista vazia quando nenhuma cor falha), o aviso não bloqueia o salvamento (o status continua \texttt{200}) e o \texttt{GET} não o devolve. Cores fora do formato \texttt{\#RRGGBB} lançam \texttt{TypeError}: não chegam às funções, pois a validação de formato as recusa antes (RN09).

\textbf{Texto sobre a cor primária.} O branco fixo não atende ao nível AA com a paleta padrão: sobre a primária \texttt{\#3b82f6}, ele dá 3,68:1. Por isso o texto sobre um fundo da cor primária (o botão principal e o item ativo da navegação) usa a variável \arq{--vitrine-on-primary}, exposta como o utilitário \texttt{text-on-primary}, cujo valor é \texttt{readableTextOn(primaryColor)}; com a paleta padrão, \texttt{\#171717}, que dá 4,87:1. O \en{layout} raiz define essa variável junto com as três cores no atributo \texttt{style} de \texttt{<html>}, e quem sobrescrever \arq{--vitrine-primary} deve sobrescrever também \arq{--vitrine-on-primary}; do contrário, uma primária escura ficaria com texto escuro. O texto sobre a cor secundária continua branco (\arq{--vitrine-on-secondary}), protegido pelo par da secundária. A decisão, tomada com o Engenheiro Chefe, consta na \adr{ADR-020}.

\textbf{Limites conhecidos.} (i) Como o texto escuro é \texttt{\#171717}, e não o preto, existe uma faixa de primárias de luminância média (luminância relativa entre 0,183 e 0,214, cerca de 6\% das cores \texttt{\#RRGGBB}; por exemplo, \texttt{\#7b7b7b} e \texttt{\#676fe3}) em que nem o branco nem o escuro chegam a 4,5:1; o pior caso é 4,23:1. Nessa faixa o par da primária também reprova e o aviso da RN09 é emitido, de modo que o administrador é informado. Com o preto puro (\texttt{\#000000}), o pior caso passa a 4,58:1 e a faixa desaparece; essa troca foi decidida na \adr{ADR-020} e está pendente na UI. (ii) A primária padrão tem 3,68:1 também sobre a página, e a de destaque padrão, 2,15:1: com a paleta padrão da RN08, texto na cor primária (links) não atende ao nível AA, e marcadores na cor de destaque ficam abaixo de 3:1; o primeiro \texttt{PUT /api/configuration} com essa paleta devolve, portanto, dois avisos (primária e destaque). Por isso os links do conteúdo (M07 e M08) não podem depender apenas da cor primária para atingir 4,5:1 com uma paleta qualquer.

\vitrinependencia{UI}{Trocar o texto escuro de \texttt{readableTextOn}, e o valor padrão de \arq{--vitrine-on-primary}, de \texttt{\#171717} para \texttt{\#000000}, ajustando \arq{tests/ui/theme.test.js}; e propor, no pedido do M07 ou do M08, o mecanismo de cor dos links do conteúdo, que atenda a 4,5:1 com uma paleta qualquer.}

\vitrinefigura[tipo=diagrama, arquivo=seq-personalizacao-tema, agente=Arquiteto, largura=0.95,
  origem={doc/diagrams/src/seq-personalizacao-tema.puml},
  descricao={Sequência de personalização: formulário, PUT na API, validação de hexadecimais e contraste, atualização no banco, revalidação do layout e leitura da configuração pelo RootLayout na próxima visita.}]
  {Sequência de personalização visual}{fig:seq-personalizacao-tema}

\section{Interface de Programação HTTP}\label{sec:arq-api}

\subsection{Convenções}

As rotas JSON seguem um REST pragmático (Seção~\ref{sec:fund-renderizacao}). Toda rota de conteúdo e de configuração exige sessão válida, verificada no servidor antes de qualquer outra ação. Entradas e saídas são JSON em UTF-8; as rotas de escrita aceitam apenas \texttt{Content-Type: application/json}, e todas as respostas levam \texttt{Cache-Control: no-store}. Códigos de status: \texttt{200}/\texttt{201}/\texttt{204} para sucesso; \texttt{400} para JSON malformado ou \texttt{Content-Type} diferente de JSON (código \arq{VALIDATION_ERROR}); \texttt{401} sem sessão válida ou com mutação vinda de outra origem; \texttt{404} recurso inexistente; \texttt{409} conflito (por exemplo, \en{slug} duplicado); \texttt{422} falha de validação; \texttt{429} excesso de tentativas de login (resposta do better-auth); \texttt{500} erro interno. Essas convenções foram aprovadas na \adr{ADR-018}.

\textbf{Resposta de sucesso.} Toda resposta de sucesso com corpo usa o envelope \texttt{\{ "data": ... \}}, que deixa espaço para campos adicionais ao lado de \texttt{data}, como \texttt{warnings} (Seção~\ref{sec:arq-contraste}). O \texttt{POST} de criação responde \texttt{201} e o cabeçalho \texttt{Location} com o caminho do novo recurso; o \texttt{DELETE} responde \texttt{204} sem corpo; e a reordenação devolve em \texttt{data} a árvore atualizada.

\begin{lstlisting}[language=VJSON, caption={Formato padrão de sucesso}, label={lst:sucesso-json}]
{
  "data": { "id": "cm...", "title": "Introdução", "slug": "introducao", "status": "PUBLISHED" }
}
\end{lstlisting}

\textbf{Resposta de erro.} Os erros têm formato único; \texttt{details} é uma lista com um item por campo, no formato \texttt{\{ field, message \}} (validação) ou \texttt{\{ field, value \}} (valor recusado). Um \texttt{409} de \en{slug} traz em \texttt{details[0].suggestion} um \en{slug} livre entre os irmãos de destino (UC05, fluxo A1).

\begin{lstlisting}[language=VJSON, caption={Formato padrão de erro}, label={lst:erro-json}]
{
  "error": {
    "code": "SLUG_CONFLICT",
    "message": "Já existe um conteúdo com este slug na mesma coleção.",
    "details": [{ "field": "slug", "value": "introducao", "suggestion": "introducao-2" }]
  }
}
\end{lstlisting}

Códigos de erro previstos: \arq{VALIDATION_ERROR}, \arq{UNAUTHENTICATED}, \arq{NOT_FOUND}, \arq{SLUG_CONFLICT}, \arq{INVALID_PARENT}, \arq{CYCLE_DETECTED} e \arq{INTERNAL_ERROR}.

\textbf{Regras comuns de entrada.} Campos desconhecidos são descartados, inclusive \texttt{authorId}, que vem sempre da sessão (RN13). Um \texttt{PATCH} ou \texttt{PUT} sem nenhum campo conhecido responde \texttt{422}. O \en{slug} reservado também é conferido ao mover um conteúdo para a raiz (por \texttt{PATCH} ou em lote). A mutação cujo cabeçalho \texttt{Origin} indica outro site é recusada com \texttt{401} (\arq{UNAUTHENTICATED}), pois o cookie de sessão não é aceito em requisições de outra origem e o contrato não prevê \texttt{403} (Seção~\ref{sec:arq-seguranca}).

\subsubsection{Erros da Camada de Dados}\label{sec:arq-dataerror}

As funções de \arq{src/lib/db.js} lançam a classe \texttt{DataError} (com os campos \texttt{code}, \texttt{message} e \texttt{details}) quando uma regra que depende do estado do banco é violada: existência, unicidade, tipo do pai, ciclo e domínio de \texttt{type}, \texttt{status} e \texttt{order}. A validação de \emph{formato} (expressão regular e \en{slugs} reservados da RN03, URL \texttt{http}/\texttt{https} da RN17, cor \texttt{\#RRGGBB} da RN09) permanece nos validadores do Core (\arq{validators.js}). Cada Route Handler captura o \texttt{DataError} e o traduz conforme a Tabela~\ref{tab:dataerror}; qualquer outro erro vira \texttt{500} com o código \arq{INTERNAL_ERROR} e mensagem genérica, sem expor detalhes internos. As funções de escrita não chamam \arq{revalidatePath} (RN10): a revalidação é da camada de rota.

\begingroup\small
\begin{longtable}{@{}L{3.9cm}L{1.2cm}L{3.3cm}L{5.8cm}@{}}
\caption{Mapeamento dos códigos de \texttt{DataError} para a resposta HTTP}\label{tab:dataerror}\\
\toprule
\textbf{Código do \texttt{DataError}} & \textbf{HTTP} & \textbf{Código da API} & \textbf{Situação}\\
\midrule
\endfirsthead
\toprule
\textbf{Código do \texttt{DataError}} & \textbf{HTTP} & \textbf{Código da API} & \textbf{Situação}\\
\midrule
\endhead
\bottomrule
\endfoot
\arq{NOT_FOUND} & 404 & \arq{NOT_FOUND} & Conteúdo (ou, ao reordenar, algum item do lote) inexistente\\
\addlinespace
\arq{SLUG_CONFLICT} & 409 & \arq{SLUG_CONFLICT} & \en{Slug} repetido entre irmãos ou na raiz (RN03)\\
\addlinespace
\arq{PARENT_NOT_FOUND} & 422 & \arq{INVALID_PARENT} & O pai informado não existe\\
\addlinespace
\arq{PARENT_NOT_COLLECTION} & 422 & \arq{INVALID_PARENT} & O pai informado é um documento (RN02)\\
\addlinespace
\arq{CYCLE} & 422 & \arq{CYCLE_DETECTED} & O destino é o próprio nó ou um descendente (RN06)\\
\addlinespace
\arq{TYPE_HAS_CHILDREN} & 422 & \arq{VALIDATION_ERROR} & Coleção com filhos não pode virar documento (\texttt{details.field} = \texttt{type})\\
\addlinespace
\arq{AUTHOR_NOT_FOUND} & 422 & \arq{VALIDATION_ERROR} & Autor inexistente; como o autor vem da sessão, indica sessão inconsistente\\
\addlinespace
\arq{INVALID_INPUT} & 422 & \arq{VALIDATION_ERROR} & Campo obrigatório ausente ou fora do domínio (\texttt{type}, \texttt{status}, \texttt{order}, lote de reordenação)\\
\addlinespace
(outro erro) & 500 & \arq{INTERNAL_ERROR} & Falha não prevista; registrar no servidor, resposta genérica\\
\end{longtable}
\fonte{Elaborado pelos autores, a partir dos códigos de \arq{src/lib/db.js} e da sugestão do agente Dados (merge M-003).}
\endgroup

Duas observações de implementação. Em \texttt{reorderContents}, os itens do lote são aplicados em sequência e validados contra o estado resultante dos anteriores, de modo que a ordem dos itens importa quando há ciclo intermediário. E o redirecionamento de uma coleção ao primeiro documento publicado (RN14) não tem função própria: o Core usa \texttt{getContentTree} com \texttt{onlyPublished} e \texttt{resolvePath}.

\subsection{Contrato dos \emph{Endpoints}}

\begin{longtable}{@{}L{1.5cm}L{3.9cm}L{1.6cm}L{4.4cm}L{2.4cm}@{}}
\caption{Contrato das rotas HTTP (normativo)}\label{tab:api}\\
\toprule
\textbf{Método} & \textbf{Caminho} & \textbf{Sessão} & \textbf{Descrição} & \textbf{Sucesso}\\
\midrule
\endfirsthead
\toprule
\textbf{Método} & \textbf{Caminho} & \textbf{Sessão} & \textbf{Descrição} & \textbf{Sucesso}\\
\midrule
\endhead
\bottomrule
\endfoot
\texttt{GET} & \arq{/api/contents} & Sim & Árvore completa, inclusive rascunhos, sem o corpo dos documentos e com o \texttt{status} de cada nó & 200\\[1pt]
\texttt{POST} & \arq{/api/contents} & Sim & Cria conteúdo (\texttt{title}, \texttt{slug}, \texttt{type}, \texttt{parentId}, \texttt{body}, \texttt{order}, \texttt{status}; este, se omitido, assume \texttt{PUBLISHED}); o autor vem da sessão & 201 e \texttt{Location}\\[1pt]
\texttt{GET} & \arq{/api/contents/:id} & Sim & Detalhe, incluindo o corpo e o \texttt{status} & 200\\[1pt]
\texttt{PATCH} & \arq{/api/contents/:id} & Sim & Atualização parcial, inclusive de \texttt{status}; dispara revalidação (RN10, RN16) & 200\\[1pt]
\texttt{DELETE} & \arq{/api/contents/:id} & Sim & Exclui o nó e seus descendentes (RN07) & 204\\[1pt]
\texttt{PATCH} & \arq{/api/contents/reorder} & Sim & Lote \texttt{[\{id, parentId, order\}]}, em transação (RN05, RN06); devolve a árvore atualizada & 200\\[1pt]
\texttt{GET} & \arq{/api/configuration} & Sim & Configuração vigente & 200\\[1pt]
\texttt{PUT} & \arq{/api/configuration} & Sim & Atualização \emph{parcial} (campos ausentes não mudam) de identidade, cores e dos seis links sociais (RN17); logotipo e \en{favicon} aceitam URL \texttt{http(s)} ou caminho de \arq{public/} iniciado por uma única barra (D-10); devolve \texttt{warnings} de contraste & 200\\[1pt]
\texttt{POST} & \arq{/api/auth/sign-in/email} & Não & Login (rota do better-auth), com limite de 5 tentativas por minuto por IP & 200 (429 acima do limite)\\[1pt]
\texttt{POST} & \arq{/api/auth/sign-up/email} & Não & Recusado: cadastro público desabilitado (RF22) & 400 (\arq{EMAIL_PASSWORD_SIGN_UP_DISABLED})\\[1pt]
\texttt{POST} & \arq{/api/auth/sign-out} & Sim & Encerra a sessão (RF20); o painel usa, em seu lugar, a \en{Server Function} \texttt{signOut} (\arq{admin/_lib/actions.js}) & 200\\[1pt]
\texttt{GET} & \arq{/api/auth/get-session} & Não & Sessão corrente ou nula & 200\\[1pt]
\end{longtable}
\fonte{Elaborado pelos autores. As rotas \arq{/api/auth/*} são as padrão do better-auth, conferidas na versão 1.7.6 (merge M-004); as demais rotas padrão da biblioteca (por exemplo, troca de senha) também ficam expostas em \arq{/api/auth/*} e exigem sessão.}

A leitura pública da árvore, dos documentos e da configuração \textbf{não} usa essas rotas: ocorre em Server Components, por chamadas diretas às funções de \arq{db.js}, o que mantém as páginas estáticas e aplica o filtro de publicação (RN15). As rotas de conteúdo e de configuração servem ao painel administrativo.

Os campos de entrada têm as seguintes regras de validação no servidor (\texttt{422} em caso de falha): \texttt{status} deve ser \texttt{DRAFT} ou \texttt{PUBLISHED}; cada link social, quando preenchido, deve ser uma URL absoluta com protocolo \texttt{http} ou \texttt{https}, e a cadeia vazia equivale a \texttt{null} (RN17); as cores devem ser \texttt{\#RRGGBB} (RN09); o \en{slug} segue a RN03 (1 a 80 caracteres, sem \en{slugs} reservados na raiz). A resposta de \texttt{422} traz um item de \texttt{details} por campo inválido.

\section{Roteamento e Navegação}\label{sec:arq-rotas}

A Figura~\ref{fig:nav-mapa-do-site} resume as rotas. Segmentos estáticos (\texttt{admin}, \texttt{login}, \texttt{api}) têm precedência sobre o \en{catch-all} \arq{[...slug]}, razão pela qual esses nomes são reservados no nível raiz (RN03). A estrutura visual de cada página segue os protótipos da IC (Seção~\ref{sec:prototipagem}): barra superior com logotipo, nome e coleções raiz; \en{sidebar} da coleção corrente; área de conteúdo; e rodapé.

\vitrinefigura[tipo=diagrama, arquivo=nav-mapa-do-site, agente=Arquiteto, largura=0.9,
  origem={doc/diagrams/src/nav-mapa-do-site.puml},
  descricao={Mapa mental das rotas: área pública (início e leitura de documento), login, painel administrativo (visão geral, árvore, novo conteúdo, editor, configurações) e API.}]
  {Mapa de rotas e navegação}{fig:nav-mapa-do-site}

\begingroup\small
\begin{longtable}{@{}L{3.3cm}L{2.9cm}L{6.0cm}L{2.2cm}@{}}
\caption{Inventário de componentes de interface e situação após o merge M-005}\label{tab:componentes-ui}\\
\toprule
\textbf{Componente} & \textbf{Diretório} & \textbf{Responsabilidade} & \textbf{Situação}\\
\midrule
\endfirsthead
\toprule
\textbf{Componente} & \textbf{Diretório} & \textbf{Responsabilidade} & \textbf{Situação}\\
\midrule
\endhead
\bottomrule
\endfoot
Navbar & \arq{components/layout} & Logotipo, nome do site e coleções raiz (RF13) & Planejado\\
Sidebar & \arq{components/layout} & Árvore da coleção corrente, item ativo, recolher níveis (RF12) & Planejado\\
Footer, SocialLinks & \arq{components/layout} & Rodapé: direitos gerados a partir do nome do site e do ano, e ícones apenas das redes com link preenchido, com os atributos \texttt{noopener} e \texttt{noreferrer} (RF24, RN17) & Planejado\\
MarkdownRenderer & \arq{components/markdown} & Pipeline de Markdown, matemática e código (RF09 a RF11) & Planejado\\
Editor, Toolbar, Preview & \arq{components/editor} & Edição com painel dividido, rascunho local e publicação (RF08, RF15, RF16) & Planejado\\
AdminShell & \arq{components/admin} & Casca do painel: atalho para o conteúdo, barra superior, navegação e área principal; \en{Server Component} de apresentação, que não valida a sessão (M10) & Implementado\\
AdminHeader & \arq{components/admin} & Barra superior: nome do site, nome e e-mail do usuário e botão ``Sair'', que recebe a \en{Server Function} \texttt{signOut} por \en{prop} (RF20) & Implementado\\
AdminSidebar & \arq{components/admin} & Navegação do painel, com item ativo pela rota e recolhível em telas estreitas (RNF06); Componente Cliente; seções sem página aparecem como ``em breve'' & Implementado\\
AdminOverview & \arq{components/admin} & Visão geral provisória: identifica a sessão e anuncia as seções futuras & Implementado\\
ContentTree & \arq{components/admin} & Árvore do painel com criar, editar, mover e excluir, mostrando o estado de cada nó (RF03, RF06, RF07, RF26) & Planejado\\
ContentForm & \arq{components/admin} & Formulário de título, \en{slug}, pai e estado de publicação (RF05, RF26) & Planejado\\
StatusToggle & \arq{components/admin} & Indicador e alternância \texttt{DRAFT}/\texttt{PUBLISHED} na árvore e no editor, com destaque para rascunhos (RF26, RN16) & Planejado\\
ThemeForm & \arq{components/admin} & Identidade, cores e links sociais, com pré-visualização e aviso de contraste (RF18, RF19, RF24) & Planejado\\
LoginForm & \arq{components/admin} & Formulário de autenticação (RF01, UC08): mensagem de erro única, aviso de excesso de tentativas e falha de conexão; Componente Cliente & Implementado\\
Button, icons & \arq{components/ui} & Botão com as variantes \texttt{primary} e \texttt{inverse}, só com tokens de tema, e ícones em SVG embutido, decorativos (RF19) & Implementado\\
Demais primitivos & \arq{components/ui} & Diálogo de confirmação, \en{toast}, campo de cor & Planejado\\
\end{longtable}
\fonte{Elaborado pelos autores. O inventário nasceu como proposta para o agente UI; os componentes ``Implementado'' vieram do merge M-005.}
\endgroup

\vitrinefigura[tipo=diagrama, arquivo=seq-reordenacao-arvore, agente=Arquiteto, largura=0.95,
  origem={doc/diagrams/src/seq-reordenacao-arvore.puml},
  descricao={Sequência de reordenação e movimentação na árvore: atualização otimista, PATCH em lote, validação de pai, ciclos e slugs, transação e revalidação do layout.}]
  {Sequência de reordenação e movimentação de conteúdos}{fig:seq-reordenacao-arvore}

\section{Segurança}\label{sec:arq-seguranca}

A Tabela~\ref{tab:ameacas} relaciona as ameaças consideradas aos controles previstos, tomando como referência o OWASP Top 10 \cite{owasptop10} e as práticas de verificação do ASVS \cite{owaspasvs}.

\begin{longtable}{@{}L{3.3cm}L{4.4cm}L{5.4cm}@{}}
\caption{Ameaças e controles}\label{tab:ameacas}\\
\toprule
\textbf{Ameaça} & \textbf{Vetor} & \textbf{Controle}\\
\midrule
\endfirsthead
\toprule
\textbf{Ameaça} & \textbf{Vetor} & \textbf{Controle}\\
\midrule
\endhead
\bottomrule
\endfoot
XSS armazenado & Markdown com HTML ou URL \texttt{javascript:} & Sem HTML bruto (RN11); manter a transformação de URL do \arq{react-markdown}; revisar qualquer \en{plugin} novo\\[1pt]
CSRF & Requisição de outro site com o \en{cookie} do administrador & \en{Cookie} \texttt{SameSite=Lax}; verificação de origem nas mutações (um \texttt{Origin} diferente do próprio site, ou \texttt{null}, responde \texttt{401} \arq{UNAUTHENTICATED} e nada é gravado; \adr{ADR-018}); rotas de escrita aceitam apenas JSON\\[1pt]
Injeção de SQL & Entradas em consultas & Prisma parametriza; proibido \texttt{\$queryRawUnsafe}\\[1pt]
Força bruta e enumeração & Tentativas repetidas de login & No máximo cinco tentativas por minuto por IP (resposta 429); mensagem genérica de erro, igual para e-mail inexistente e senha errada (UC08, CT02); o \en{proxy} deve definir \texttt{x-forwarded-for} (Seção~\ref{sec:arq-auth})\\[1pt]
Acesso indevido à API & Chamada direta sem passar pelo painel & Sessão verificada em cada Route Handler (RF21), antes de ler o corpo; Proxy apenas otimista\\[1pt]
Rotas padrão do better-auth & Rotas de \arq{/api/auth/*} além do login (por exemplo, troca de senha e listagem de sessões) & Exigem sessão e passam pela mesma limitação de requisições; o cadastro público está desabilitado (RF22)\\[1pt]
Exposição de segredos & \arq{.env} no repositório, na imagem Docker ou variáveis no navegador & \arq{.env} ignorado pelo Git e excluído da imagem por \arq{.dockerignore}, fornecido ao contêiner em tempo de execução; nenhuma variável com prefixo \arq{NEXT_PUBLIC_} para segredos\\[1pt]
Vazamento de rascunho & Nó \texttt{DRAFT} acessível por URL, pela árvore ou por página em cache & Filtro de publicação na camada de dados (RN15); revalidação do \en{layout} ao alterar o estado (RN16); testes CT31 e CT32\\[1pt]
Link social malicioso & URL \texttt{javascript:} ou \texttt{data:} em um dos seis campos do rodapé & Validação de protocolo \texttt{http}/\texttt{https} no servidor (RN17); atributos \texttt{noopener} e \texttt{noreferrer}; teste CT33\\[1pt]
Injeção de CSS & Valor de cor malicioso & Validação \texttt{\#RRGGBB} no servidor (RN09)\\[1pt]
Dependências vulneráveis & Pacotes desatualizados, com vulnerabilidade conhecida ou maliciosos & Revisão de cada nova dependência pelo Arquiteto; \texttt{npm audit} a cada entrega; \en{lockfile} versionado; os quatro alertas atuais foram avaliados e aceitos (Seção~\ref{sec:arq-npm-audit}, \adr{ADR-019})\\[1pt]
\en{Clickjacking} & Página incorporada em \texttt{iframe} & Cabeçalho \arq{frame-ancestors}/\arq{X-Frame-Options}\\[1pt]
Exclusão acidental & Cascata indevida & Confirmação explícita; \en{backup} (RNF10)\\[1pt]
Perda do banco & Remoção do diretório de dados, falha do disco da VPS ou atualização que descarta o volume & Banco em \en{bind mount} fora da imagem; \en{backup} periódico com \texttt{VACUUM INTO} e cópia externa (Apêndice~\ref{ap:implantacao}; risco R12)\\[1pt]
\end{longtable}
\fonte{Elaborado pelos autores.}

Além dos controles acima, o \arq{next.config.mjs} deve definir cabeçalhos de segurança (\texttt{X-Content-Type-Options: nosniff}, \arq{Referrer-Policy}, política de enquadramento e \arq{Permissions-Policy} restritiva); HSTS e TLS ficam a cargo do \en{proxy} reverso. Uma política CSP com \en{nonce} é desejável e será avaliada pelo agente Core.

\subsection{Vulnerabilidades Conhecidas em Dependências}\label{sec:arq-npm-audit}

Em 01/10/2026, na revisão do merge M-004, o \texttt{npm audit} acusou \textbf{quatro vulnerabilidades de severidade alta}, todas em dependências transitivas do Prisma 7 (Tabela~\ref{tab:npm-audit}). O npm só as corrige com \texttt{npm audit fix -{}-force}, que instalaria o \texttt{prisma@6.19.3}: um rebaixamento de versão maior, incompatível com o que o projeto adota (\arq{prisma.config.ts}, gerador \texttt{prisma-client} e \en{driver adapter}; \adr{ADR-014}). A \texttt{7.10.0}, instalada, é a última versão 7 publicada, e o \texttt{npm audit} lista como afetada toda a faixa do \arq{prisma} até a \texttt{8.1.0-dev.6}; portanto, não há versão 7 corrigida. Decidiu-se \textbf{manter o Prisma 7 e não aplicar a correção automática}, aceitando os alertas, pelas razões da tabela (\adr{ADR-019}). Como \arq{prisma} agora está em \texttt{dependencies} (para aplicar migrações na imagem), esses pacotes também estão presentes na imagem de produção, embora não sejam exercidos.

\begingroup\small
\begin{longtable}{@{}L{3.3cm}L{5.4cm}L{5.5cm}@{}}
\caption{Alertas do \texttt{npm audit} aceitos em 01/10/2026}\label{tab:npm-audit}\\
\toprule
\textbf{Pacote} & \textbf{Aviso} & \textbf{Avaliação para o Vitrine}\\
\midrule
\endfirsthead
\toprule
\textbf{Pacote} & \textbf{Aviso} & \textbf{Avaliação para o Vitrine}\\
\midrule
\endhead
\bottomrule
\endfoot
\arq{deepmerge-ts} (antes da 8.0.0; usado por \arq{@prisma/config}) & Esgotamento de pilha ao mesclar grafos de objetos recursivos (GHSA-ggr8-5vv4-36mx) & Roda na CLI do Prisma, ao mesclar a configuração de \arq{prisma.config.ts}, escrita pela equipe; não recebe entrada de usuário nem é exercido pelo servidor\\
\addlinespace
\arq{@prisma/config} e \arq{prisma} & Alertas em cadeia, por dependerem do \arq{deepmerge-ts} e do \arq{mysql2} & O mesmo das linhas vizinhas; sem correção em nenhuma versão 7 disponível\\
\addlinespace
\arq{mysql2} (até a 3.23.0; dependência do \arq{prisma} e do \arq{better-auth}) & Rebaixamento do plugin de autenticação para \arq{mysql_clear_password}, com vazamento de credenciais (GHSA-3f6p-5ww8-9rcr), e descompressão (\en{inflate}) sem limite no protocolo comprimido, com negação de serviço (GHSA-rgwj-5xj2-c3m3) & Só é exercido ao conectar a um servidor MySQL; o Vitrine usa SQLite por \arq{better-sqlite3} e nunca abre conexão MySQL\\
\end{longtable}
\fonte{Elaborado pelos autores, a partir da saída do \texttt{npm audit} e de \texttt{npm ls} (merge M-004).}
\endgroup

A decisão será reavaliada a cada entrega (o \texttt{npm audit} é critério de saída, Seção~\ref{sec:testes-saida}) e sempre que o Prisma publicar uma versão 7 com as dependências corrigidas. Enquanto isso, os alertas aceitos não contam como vulnerabilidade crítica pendente.

\section{Implantação e Operação}\label{sec:arq-implantacao}

Por decisão do Engenheiro Chefe (\decisao{D-04}, \adr{ADR-011}), o alvo de implantação é uma VPS tradicional com Docker Engine e Docker Compose. A aplicação é empacotada em uma imagem construída em dois estágios (instalação e \en{build}; execução) que executa o servidor Next.js (\texttt{next start}) como usuário sem privilégios, em um único contêiner. O arquivo do SQLite fica em um \en{bind mount} do hospedeiro (\arq{./data}, montado em \arq{/data}), fora da imagem, e as variáveis de ambiente vêm de um arquivo \arq{.env} também fora da imagem. Um \en{proxy} reverso com TLS (Caddy ou Nginx) encaminha as requisições para a porta 3000 do contêiner, publicada apenas em \arq{127.0.0.1} (Figura~\ref{fig:impl-implantacao}). O passo a passo, com esboços do \arq{Dockerfile} e do \arq{docker-compose.yml}, está no Apêndice~\ref{ap:implantacao}.

\vitrinefigura[tipo=diagrama, arquivo=impl-implantacao, agente=Arquiteto, largura=0.95,
  origem={doc/diagrams/src/impl-implantacao.puml},
  descricao={Diagrama de implantação: VPS com Docker Engine; proxy reverso com TLS; contêiner vitrine com o build do Next.js, o cliente do Prisma e o adaptador better-sqlite3; volume ./data com o arquivo SQLite (e WAL) e cópias de segurança; arquivo .env fora da imagem.}]
  {Diagrama de implantação}{fig:impl-implantacao}

Os requisitos normativos sobre a imagem, que o módulo M13 deve atender, são:

\begin{itemize}
  \item \textbf{Módulo nativo.} O \arq{better-sqlite3} é um módulo nativo (\adr{ADR-014}); os estágios de \en{build} e de execução devem usar a mesma versão do Node.js e a mesma distribuição (\en{glibc}, por exemplo \texttt{node:22-bookworm-slim}), e o estágio de \en{build} precisa de compilador (\texttt{python3}, \texttt{make}, \texttt{g++}) caso não haja binário pré-compilado;
  \item \textbf{Cliente gerado.} \texttt{prisma generate} roda antes de \texttt{next build}, pois \arq{src/generated/prisma} não é versionado;
  \item \textbf{Migrações como passo explícito.} \texttt{prisma migrate deploy} é executado por um comando próprio antes de subir a nova versão (por exemplo, \texttt{docker compose run -{}-rm vitrine npx prisma migrate deploy}), e não pelo comando de partida da aplicação; isso exige que a CLI do Prisma e o \arq{dotenv} estejam na imagem, o que está garantido desde o merge M-002, que os moveu para \texttt{dependencies} (\adr{ADR-014});
  \item \textbf{Administrador inicial.} Depois das migrações e do \en{seed}, o administrador é criado por \texttt{docker compose run -{}-rm vitrine npm run admin:create}, que exige \arq{ADMIN_EMAIL} e \arq{ADMIN_PASSWORD} (e aceita \arq{ADMIN_NAME}) no ambiente da execução. O \en{seed} (\texttt{prisma db seed}) \emph{não} cria o administrador: ele só garante a configuração padrão (\adr{ADR-018});
  \item \textbf{\texttt{tsx} na imagem.} O \texttt{admin:create} e o \texttt{prisma db seed} rodam com o \texttt{tsx}, hoje em \texttt{devDependencies}; a imagem de produção deve manter as dependências de desenvolvimento (\texttt{npm ci} completo) ou mover o \texttt{tsx} para \texttt{dependencies}, caso contrário esses dois comandos deixam de funcionar;
  \item \textbf{\texttt{DATABASE\_URL} no \en{build}.} O \texttt{next build} importa as rotas, que importam \arq{db.js}, e este lança erro se \arq{DATABASE_URL} não está definida; o estágio de \en{build} do Dockerfile deve defini-la com qualquer valor bem formado (por exemplo, \texttt{ENV} \arq{DATABASE_URL=file:/data/vitrine.db}), pois nenhuma consulta ocorre no \en{build} e o arquivo não precisa existir. Segredos (\arq{BETTER_AUTH_SECRET}, \arq{ADMIN_PASSWORD}) \emph{não} entram no \en{build}: o aviso do segredo padrão registrado nessa etapa é esperado e inofensivo;
  \item \textbf{Segredo em execução.} \arq{BETTER_AUTH_SECRET} e \arq{BETTER_AUTH_URL} (em \texttt{https}, para o \en{cookie} \texttt{Secure}) são obrigatórios na execução; sem o segredo, as rotas de autenticação e de API respondem 500;
  \item \textbf{Instância única.} Um único contêiner, pois o cache de ISR é por instância (Seção~\ref{sec:arq-isr}), o SQLite tem um escritor por vez e a limitação de tentativas de login guarda os contadores em memória;
  \item \textbf{Segredos.} \arq{.env} e \arq{data/} constam do \arq{.dockerignore} e do \arq{.gitignore}.
\end{itemize}

O relatório de IC citava a plataforma Vercel como exemplo de implantação \cite{reis2026ic}. O sistema de arquivos de funções \emph{serverless} é efêmero, de modo que um arquivo SQLite local não persistiria entre invocações nesse tipo de ambiente; a decisão por uma VPS com volume elimina esse obstáculo, que constava como risco R01, para esta versão. Plataformas dessa natureza só seriam viáveis com a troca do \en{driver adapter} por um banco remoto compatível (por exemplo, libSQL/Turso), o que está fora do escopo (\adr{ADR-011}).

\section{Desempenho e Escalabilidade}\label{sec:arq-desempenho}

O perfil de carga é dominado por leituras. A ISR evita consultas ao banco na maioria das requisições públicas; quando há regeneração, a consulta da árvore é linear no número de nós e a resolução do caminho é proporcional à profundidade. O SQLite em modo WAL permite leitores simultâneos à escrita, que é serializada \cite{sqlitewal}; como as escritas ocorrem apenas por ação do administrador, a limitação de um escritor por vez não é relevante para o perfil esperado \cite{sqlitewhentouse}. Os limites práticos --- tamanho da árvore, número de documentos pré-renderizados no \en{build} e tempo de regeneração do \en{layout} raiz --- serão medidos no Capítulo~\ref{cap:testes}.

\section{Síntese do Capítulo}

O capítulo especificou o monólito modular, suas visões, regras de dependência, modelo de dados (com o estado de publicação e os links sociais), autenticação em camadas, estratégia de cache e de visibilidade, pipeline de conteúdo, tematização, contrato de API, segurança e implantação por imagem Docker. As observações sobre o esquema atual (Tabela~\ref{tab:obs-schema}) e os valores padrão da \adr{ADR-017} formam a lista de tarefas de refinamento dos agentes; o Capítulo~\ref{cap:implementacao} acompanha a execução.
